\documentclass[a4paper,12pt]{article}

\usepackage[utf8]{inputenc}
\usepackage[T1]{fontenc}
\usepackage{amsmath}
\usepackage{booktabs}
\usepackage{graphicx}
\usepackage{geometry}
\geometry{margin=2.5cm}
\usepackage{xcolor}
\usepackage{listings}
\usepackage{hyperref}
\hypersetup{colorlinks=true, linkcolor=blue!50!black, urlcolor=blue!50!black}

% Nomes em português (o ambiente de compilação não possui os padrões de
% hifenização do babel para português instalados, por isso os nomes são
% redefinidos manualmente em vez de carregar o pacote babel)
\renewcommand{\tablename}{Tabela}
\renewcommand{\lstlistingname}{Código}

% ---------- Configuração do realce de sintaxe (C) ----------
\definecolor{codebg}{RGB}{247,247,247}
\definecolor{codekw}{RGB}{0,0,180}
\definecolor{codecomment}{RGB}{0,128,0}
\definecolor{codestring}{RGB}{163,21,21}
\definecolor{framecolor}{RGB}{200,200,200}

\lstdefinestyle{cstyle}{
    language=C,
    backgroundcolor=\color{codebg},
    basicstyle=\ttfamily\footnotesize,
    keywordstyle=\color{codekw}\bfseries,
    commentstyle=\color{codecomment}\itshape,
    stringstyle=\color{codestring},
    numberstyle=\tiny\color{gray},
    numbers=left,
    numbersep=6pt,
    frame=single,
    rulecolor=\color{framecolor},
    breaklines=true,
    breakatwhitespace=false,
    showstringspaces=false,
    tabsize=4,
    captionpos=b,
    columns=flexible,
    morekeywords={size_t},
    extendedchars=true,
    literate=
        {á}{{\'a}}1 {à}{{\`a}}1 {â}{{\^a}}1 {ã}{{\~a}}1
        {é}{{\'e}}1 {ê}{{\^e}}1 {í}{{\'i}}1
        {ó}{{\'o}}1 {ô}{{\^o}}1 {õ}{{\~o}}1
        {ú}{{\'u}}1 {ü}{{\"u}}1 {ç}{{\c c}}1
        {Á}{{\'A}}1 {À}{{\`A}}1 {Â}{{\^A}}1 {Ã}{{\~A}}1
        {É}{{\'E}}1 {Ê}{{\^E}}1 {Í}{{\'I}}1
        {Ó}{{\'O}}1 {Ô}{{\^O}}1 {Õ}{{\~O}}1
        {Ú}{{\'U}}1 {Ü}{{\"U}}1 {Ç}{{\c C}}1
}

\title{Multiplicação de Matriz por Vetor: \\ Localidade Temporal e Espacial no Acesso à Memória}
\author{Israel Soares de Castro Filho \\ Matrícula: 20250070780}
\date{}

\begin{document}

\maketitle

\section{Introdução}

Este relatório apresenta a implementação de um programa em linguagem C para
comparar duas formas de percorrer os elementos de uma matriz durante a
multiplicação de matriz por vetor (MxV), avaliando o impacto do padrão de
acesso à memória sobre o desempenho:

\begin{itemize}
    \item \textbf{Versão por linhas}: o laço externo percorre as linhas $i$ e o
    laço interno percorre as colunas $j$, de modo que os elementos $A[i][j]$ são
    acessados sequencialmente na memória (\emph{stride} 1).
    \item \textbf{Versão por colunas}: o laço externo percorre as colunas $j$ e o
    laço interno percorre as linhas $i$, de modo que os elementos $A[i][j]$ são
    acessados com um salto de $N$ posições a cada passo (\emph{stride} $N$), pois
    em C as matrizes são armazenadas em ordem \emph{row-major} (por linhas).
\end{itemize}

O objetivo é identificar a partir de que tamanho de matriz os tempos de execução
das duas versões passam a divergir de forma significativa e explicar essa
divergência com base no funcionamento da memória \emph{cache}.

\section{Metodologia}

O programa aloca uma matriz quadrada $N \times N$ como um único bloco contíguo de
memória e mede o tempo de execução de cada versão do algoritmo com a função
\texttt{clock\_gettime()} (relógio \texttt{CLOCK\_MONOTONIC} da biblioteca
\texttt{time.h}), que oferece alta resolução e não é afetada por ajustes do
relógio do sistema. Para cada tamanho de matriz $N$ (variando entre $64$ e
$16384$, em potências de 2), cada versão foi executada 5 vezes, sendo
registrado o menor tempo obtido, de forma a reduzir o ruído de medição causado
por outras tarefas do sistema operacional. O código completo, compilado com
\texttt{gcc -O2 -o tarefa2 tarefa2.c}, é apresentado no Código~\ref{lst:tarefa2}.

\section{Resultados}

A Tabela~\ref{tab:resultados} apresenta os tempos de execução (em segundos,
melhor caso entre 5 repetições) das duas versões do algoritmo, além da razão
entre o tempo da versão por colunas e o tempo da versão por linhas.

\begin{table}[h]
\centering
\caption{Tempos de execução da multiplicação MxV por linhas e por colunas}
\label{tab:resultados}
\begin{tabular}{rrrr}
\toprule
N & Linhas (s) & Colunas (s) & Razão (col./lin.) \\
\midrule
64      & 0,000004 & 0,000004 & $1{,}03\times$  \\
128     & 0,000018 & 0,000017 & $0{,}94\times$  \\
256     & 0,000075 & 0,000144 & $1{,}91\times$  \\
512     & 0,000222 & 0,001498 & $6{,}75\times$  \\
1.024   & 0,000938 & 0,008236 & $8{,}78\times$  \\
2.048   & 0,003564 & 0,051392 & $14{,}42\times$ \\
4.096   & 0,014300 & 0,379421 & $26{,}53\times$ \\
8.192   & 0,057248 & 3,403755 & $59{,}46\times$ \\
16.384  & 0,267308 & 8,399817 & $31{,}42\times$ \\
\bottomrule
\end{tabular}
\end{table}

\section{Análise}

\subsection{Ponto de divergência dos tempos}

Para matrizes pequenas ($N = 64$ e $N = 128$), os tempos das duas versões são
praticamente idênticos (razão próxima de $1\times$). A partir de $N = 256$ a
diferença começa a aparecer (razão de aproximadamente $1{,}9\times$) e cresce de
forma acentuada e consistente até $N = 8192$, onde a versão por colunas chega a
ser quase 60 vezes mais lenta que a versão por linhas. Pode-se dizer, portanto,
que a divergência \emph{significativa} de desempenho passa a ocorrer a partir de
$N \approx 256$--$512$, tornando-se dramática para $N \geq 2048$.

\subsection{Relação com a hierarquia de memória (cache)}

O comportamento observado é explicado diretamente pelo conceito de
\textbf{localidade espacial} e pela forma como a matriz é armazenada na memória.
Em C, uma matriz bidimensional é organizada em ordem \emph{row-major}: os
elementos de uma mesma linha ficam contíguos na memória, enquanto elementos de
uma mesma coluna estão separados por $N$ posições (uma linha inteira).

Na versão por linhas, o acesso sequencial (\emph{stride} 1) faz com que, ao
carregar um elemento, todo o bloco de memória ao seu redor (a \emph{cache
line}, tipicamente 64 bytes, ou seja, 8 valores \texttt{double}) seja
aproveitado nos acessos seguintes, gerando muitos \emph{cache hits}. Já na
versão por colunas, cada acesso salta $N$ posições (\emph{stride} $N$): para
valores de $N$ suficientemente grandes, cada acesso cai em uma linha de cache
diferente, gerando \emph{cache misses} frequentes e obrigando o processador a
buscar dados repetidamente na memória principal, muito mais lenta que a cache.

\subsection{Por que a diferença só aparece a partir de certo tamanho}

Para matrizes pequenas, o conjunto de dados inteiro cabe (ou quase cabe) nas
caches L1/L2 do processador, de modo que mesmo os acessos com \emph{stride}
grande acabam sendo atendidos pela cache na maior parte do tempo. À medida que
$N$ cresce, o tamanho da matriz ($N^2 \times 8$ bytes) ultrapassa a capacidade
das caches L1, depois L2 e eventualmente L3, e a taxa de \emph{cache misses} da
versão por colunas aumenta continuamente -- exatamente o efeito observado no
crescimento da razão entre os tempos até $N = 8192$. A leve queda da razão em
$N = 16384$ ocorre porque, nesse ponto, mesmo a versão por linhas já não cabe
inteiramente nas caches mais rápidas, de modo que seu tempo também passa a ser
dominado por acessos à memória principal, reduzindo a vantagem relativa entre as
duas versões.

\subsection{Localidade temporal}

Além da localidade espacial, há também um componente de localidade
\textbf{temporal} relevante: o vetor \texttt{v} é reutilizado em todas as
iterações do laço. Na versão por linhas, cada elemento \texttt{v[j]} é acessado
repetidamente em iterações próximas, favorecendo sua permanência na cache, e a
variável acumuladora \texttt{soma} é mantida em registrador durante toda a soma
de uma linha. Esses fatores reforçam ainda mais a vantagem da versão por linhas
sobre a versão por colunas.

\section{Reflexão: aplicações reais}

O padrão observado neste experimento -- programas com a mesma complexidade
assintótica apresentando desempenho drasticamente diferente conforme o padrão
de acesso à memória -- se repete em diversas áreas da computação:

\begin{itemize}
    \item \textbf{Álgebra linear e computação científica:} bibliotecas de alto
    desempenho (como BLAS/LAPACK) reorganizam laços e utilizam técnicas de
    \emph{blocking}/\emph{tiling} exatamente para maximizar a reutilização de
    dados na cache em operações com matrizes grandes, obtendo ganhos de
    desempenho semelhantes aos observados aqui.

    \item \textbf{Bancos de dados:} o formato de armazenamento de tabelas
    (orientado a linhas ou a colunas) é escolhido de acordo com o padrão de
    acesso predominante da aplicação -- sistemas OLTP tendem a usar
    armazenamento por linhas, enquanto sistemas analíticos (OLAP) costumam usar
    armazenamento por colunas, cada um otimizado para seu próprio padrão de
    varredura.

    \item \textbf{Processamento de imagens e aprendizado de máquina:} operações
    sobre tensores multidimensionais (convoluções, multiplicações de matrizes em
    redes neurais) são altamente sensíveis à ordem de acesso à memória; por
    isso, frameworks como cuDNN e cuBLAS investem pesadamente em otimizações de
    localidade de memória
\end{itemize}

\newpage
\appendix

\section{Código-Fonte}

Código utilizado para gerar os resultados apresentados neste relatório.

\lstinputlisting[style=cstyle, caption={Código-fonte \texttt{tarefa2.c} -- multiplicação matriz-vetor com dois padrões de acesso à memória.}, label={lst:tarefa2}]{../src/tarefa2.c}

\end{document}